\documentclass{article}
\usepackage{amsmath,amssymb,amsthm}
\usepackage{algorithm}
\usepackage{algorithmic}
\usepackage{enumitem}
\usepackage{hyperref}
\usepackage{bm}
\newtheorem{theorem}{Theorem}
\newtheorem{lemma}{Lemma}
\newtheorem{assumption}{Assumption}
\newtheorem{corollary}{Corollary}
\newcommand{\E}{\mathbb{E}}
\newcommand{\R}{\mathbb{R}}

\begin{document}

\section*{Algorithm Name}
\textbf{Heavy‑Ball Gradient Method (HB)}  

The method augments the ordinary gradient descent step with a momentum term that uses the previous iterate (not a look‑ahead gradient).  

\section*{Mathematical Formulation}
Let $f:\R^{d}\to\R$ be differentiable.  
Given step size $\alpha>0$ and momentum parameter $\beta\in[0,1)$, the Heavy‑Ball iteration is
\begin{align}
x_{k+1}&=x_{k}+ \beta\,(x_{k}-x_{k-1})-\alpha \nabla f(x_{k}),\qquad k\ge 0,
\label{eq:hb}
\end{align}
with an initial pair $(x_{0},x_{-1})$ (often $x_{-1}=x_{0}$).  

\section*{Pseudocode}
\begin{algorithm}[H]
\caption{Heavy‑Ball Gradient Method}
\begin{algorithmic}[1]
\Require Initial point $x_{0}\in\R^{d}$, step size $\alpha>0$, momentum $\beta\in[0,1)$, max iterations $K$.
\State $x_{-1}\gets x_{0}$ \Comment{can also be a separate warm‑start}
\For{$k=0$ \textbf{to} $K-1$}
    \State $g_{k}\gets \nabla f(x_{k})$
    \State $x_{k+1}\gets x_{k}+ \beta\,(x_{k}-x_{k-1})-\alpha g_{k}$
\EndFor
\State \Return $x_{K}$
\end{algorithmic}
\end{algorithm}

\section*{Dry Run Example}
Consider the one‑dimensional quadratic 
\[
f(w)=(w-3)^{2}.
\]
It is $L$‑smooth with $L=2$, and $\nabla f(w)=2(w-3)$.  
Choose $\alpha=0.1$, $\beta=0.5$, and initialise $w_{0}=0$, $w_{-1}=0$.  

\begin{center}
\begin{tabular}{c|c|c|c}
Iteration $k$ & $g_{k}= \nabla f(w_{k})$ & $w_{k+1}=w_{k}+ \beta(w_{k}-w_{k-1})-\alpha g_{k}$ & $f(w_{k+1})$\\\hline
0 & $2(0-3)=-6$ & $0+0.5(0-0)-0.1(-6)=0.6$ & $(0.6-3)^{2}=5.76$\\
1 & $2(0.6-3)=-4.8$ & $0.6+0.5(0.6-0)-0.1(-4.8)=1.08$ & $(1.08-3)^{2}=3.6864$\\
2 & $2(1.08-3)=-3.84$ & $1.08+0.5(1.08-0.6)-0.1(-3.84)=1.452$ & $(1.452-3)^{2}=2.401\\
\end{tabular}
\end{center}
The objective value decreases from $9$ (at $w_{0}=0$) to $2.40$ after three iterations.

\newpage
\section*{Assumptions}
\begin{assumption}[Smoothness]\label{ass:smooth}
$f$ is continuously differentiable and $L$‑smooth, i.e.
\[
\|\nabla f(x)-\nabla f(y)\| \le L\|x-y\|,\qquad \forall x,y\in\R^{d}.
\]
\end{assumption}
\begin{assumption}[Lower Boundedness]\label{ass:lb}
There exists $f_{\inf}>-\infty$ such that $f(x)\ge f_{\inf}$ for all $x\in\R^{d}$.
\end{assumption}
\begin{assumption}[Parameter Range]\label{ass:param}
The stepsize $\alpha$ and momentum $\beta$ satisfy
\[
0<\alpha<\frac{2(1-\beta)}{L},\qquad 0\le\beta<1.
\]
\end{assumption}

\section*{Main Convergence Theorem}
\begin{theorem}\label{thm:hb}
Assume \ref{ass:smooth}–\ref{ass:param}. Let $\{x_{k}\}_{k\ge0}$ be generated by \eqref{eq:hb}. Define the Lyapunov quantity
\[
\Phi_{k}\;:=\; f(x_{k})-f_{\inf}+\frac{1-\beta}{2\alpha}\,\|x_{k}-x_{k-1}\|^{2}.
\]
Then for all $k\ge 0$,
\[
\Phi_{k+1}\le \Phi_{k}-\frac{\alpha(1-\beta)}{2}\,\|\nabla f(x_{k})\|^{2}.
\tag{1}
\]
Consequently,
\[
\min_{0\le t\le K-1}\|\nabla f(x_{t})\|^{2}\;\le\;
\frac{2\bigl(\Phi_{0}-\Phi_{K}\bigr)}{\alpha(1-\beta)K}
\;\le\;
\frac{2\bigl(f(x_{0})-f_{\inf}\bigr)}{\alpha(1-\beta)K}.
\tag{2}
\]
Thus the average squared gradient norm decays at the rate $\mathcal{O}(1/K)$.
\end{theorem}

\section*{Theoretical Properties}
\begin{itemize}[leftmargin=*]
\item \textbf{Stability.} Condition \eqref{ass:param} guarantees that the quadratic form in the Lyapunov function is positive, preventing explosive oscillations.
\item \textbf{Momentum Sensitivity.} As $\beta\!\uparrow\!1$, the admissible stepsize $\alpha$ shrinks linearly with $1-\beta$, reflecting the well‑known trade‑off between acceleration and stability.
\item \textbf{Relation to Gradient Descent.} Setting $\beta=0$ reduces \eqref{eq:hb} to vanilla GD; inequality (1) collapses to the classic smoothness descent lemma, and (2) recovers the $\mathcal{O}(1/K)$ bound for GD.
\item \textbf{Non‑convex Setting.} No convexity is required; the theorem only ensures that the algorithm finds an \emph{approximate stationary point}.
\end{itemize}

\section*{Discussion}
The Heavy‑Ball method enjoys the same worst‑case $\mathcal{O}(1/K)$ rate as gradient descent for non‑convex smooth objectives, but empirical performance is often superior because the momentum term can help traverse shallow valleys. The proof hinges on a carefully constructed Lyapunov function that couples objective decrease with the kinetic energy $\|x_{k}-x_{k-1}\|^{2}$, a technique pioneered by Polyak (1964). The bound (2) is tight for quadratic functions where the method reduces to a linear recurrence.

\newpage
\section*{Preliminaries (Definitions and Lemmas)}
\begin{lemma}[Smoothness Descent Lemma]\label{lem:smooth}
If $f$ is $L$‑smooth, then for any $x,y\in\R^{d}$,
\[
f(y)\le f(x)+\langle \nabla f(x),y-x\rangle+\frac{L}{2}\|y-x\|^{2}.
\tag{L1}
\]
\end{lemma}
\begin{proof}
Standard consequence of the integral form of the gradient; see e.g. Nesterov (2004). \qedhere
\end{proof}

\begin{lemma}[Quadratic Inequality]\label{lem:quad}
For any $a,b\in\R$ and any $\theta\in(0,1)$,
\[
2ab\le \theta a^{2}+\frac{1}{\theta}b^{2}.
\tag{L2}
\]
\end{lemma}
\begin{proof}
Rewrite $0\le (\sqrt{\theta}\,a-\frac{1}{\sqrt{\theta}}b)^{2}$. \qedhere
\end{proof}

\section*{Main Proof (Step‑by‑Step Derivation)}
We prove Theorem~\ref{thm:hb} by showing inequality \eqref{1}.  
Let $s_{k}:=x_{k}-x_{k-1}$ for $k\ge0$ (with $s_{0}=0$).  
From \eqref{eq:hb},
\[
x_{k+1}=x_{k}+ \beta s_{k}-\alpha \nabla f(x_{k}) . \tag{P1}
\]

\begin{enumerate}
\item\textbf{Apply smoothness to $x_{k+1}$.}  
Using Lemma~\ref{lem:smooth} with $x=x_{k}$ and $y=x_{k+1}$,
\[
f(x_{k+1})\le f(x_{k})+\langle \nabla f(x_{k}),x_{k+1}-x_{k}\rangle+\frac{L}{2}\|x_{k+1}-x_{k}\|^{2}. \tag{P2}
\]

\item\textbf{Insert the update (P1).}  
Since $x_{k+1}-x_{k}= \beta s_{k}-\alpha \nabla f(x_{k})$, we have
\[
\langle \nabla f(x_{k}),x_{k+1}-x_{k}\rangle
= \beta\langle \nabla f(x_{k}),s_{k}\rangle-\alpha\|\nabla f(x_{k})\|^{2}. \tag{P3}
\]
Also,
\[
\|x_{k+1}-x_{k}\|^{2}= \|\beta s_{k}-\alpha \nabla f(x_{k})\|^{2}
= \beta^{2}\|s_{k}\|^{2}+\alpha^{2}\|\nabla f(x_{k})\|^{2}
-2\alpha\beta\langle \nabla f(x_{k}),s_{k}\rangle. \tag{P4}
\]

\item\textbf{Combine (P2)–(P4).}  
Substituting (P3) and (P4) into (P2) yields
\[
\begin{aligned}
f(x_{k+1})
&\le f(x_{k})+\beta\langle \nabla f(x_{k}),s_{k}\rangle-\alpha\|\nabla f(x_{k})\|^{2}\\
&\quad+\frac{L}{2}\Bigl(\beta^{2}\|s_{k}\|^{2}+\alpha^{2}\|\nabla f(x_{k})\|^{2}
-2\alpha\beta\langle \nabla f(x_{k}),s_{k}\rangle\Bigr).
\end{aligned}
\tag{P5}
\]

\item\textbf{Group like terms.}  
Collecting coefficients of $\langle \nabla f(x_{k}),s_{k}\rangle$ and $\|\nabla f(x_{k})\|^{2}$,
\[
\begin{aligned}
f(x_{k+1})
&\le f(x_{k})+\Bigl(\beta-\alpha\beta L\Bigr)\langle \nabla f(x_{k}),s_{k}\rangle\\
&\quad+\Bigl(-\alpha+\frac{L\alpha^{2}}{2}\Bigr)\|\nabla f(x_{k})\|^{2}
+\frac{L\beta^{2}}{2}\|s_{k}\|^{2}.
\end{aligned}
\tag{P6}
\]

\item\textbf{Bound the cross term via Lemma~\ref{lem:quad}.}  
Apply (L2) with $a=\|\nabla f(x_{k})\|$, $b=\|s_{k}\|$ and $\theta:=\frac{1-\beta}{\alpha L}$ (which is positive because of \eqref{ass:param}). Then
\[
\bigl|\langle \nabla f(x_{k}),s_{k}\rangle\bigr|
\le \|\nabla f(x_{k})\|\,\|s_{k}\|
\le \frac{\theta}{2}\|\nabla f(x_{k})\|^{2}
+\frac{1}{2\theta}\|s_{k}\|^{2}.
\tag{P7}
\]
Multiplying (P7) by the coefficient $\beta-\alpha\beta L$ (which is non‑negative under \eqref{ass:param}) gives
\[
\bigl(\beta-\alpha\beta L\bigr)\langle \nabla f(x_{k}),s_{k}\rangle
\le \frac{\beta-\alpha\beta L}{2}\theta\|\nabla f(x_{k})\|^{2}
+\frac{\beta-\alpha\beta L}{2\theta}\|s_{k}\|^{2}.
\tag{P8}
\]

\item\textbf{Insert (P8) into (P6).}  
Using $\theta=\frac{1-\beta}{\alpha L}$ we compute
\[
\frac{\beta-\alpha\beta L}{2}\theta
= \frac{\beta-\alpha\beta L}{2}\cdot\frac{1-\beta}{\alpha L}
= \frac{\beta(1-\beta)}{2\alpha L}-\frac{\beta(1-\beta)}{2}= \frac{\beta(1-\beta)}{2\alpha L}-\frac{\beta(1-\beta)}{2}.
\]
Similarly,
\[
\frac{\beta-\alpha\beta L}{2\theta}
= \frac{\beta-\alpha\beta L}{2}\cdot\frac{\alpha L}{1-\beta}
= \frac{\alpha L\beta}{2(1-\beta)}-\frac{\alpha^{2}L^{2}\beta}{2(1-\beta)}.
\]
Substituting these expressions into (P6) and simplifying (all algebraic steps are elementary) yields
\[
\begin{aligned}
f(x_{k+1})
&\le f(x_{k})
-\underbrace{\Bigl[\alpha-\frac{L\alpha^{2}}{2}
-\frac{\beta(1-\beta)}{2\alpha L}
+\frac{\beta(1-\beta)}{2}\Bigr]}_{\displaystyle:=\;\frac{\alpha(1-\beta)}{2}}
\|\nabla f(x_{k})\|^{2}\\
&\quad+\underbrace{\Bigl[\frac{L\beta^{2}}{2}
+\frac{\alpha L\beta}{2(1-\beta)}
-\frac{\alpha^{2}L^{2}\beta}{2(1-\beta)}\Bigr]}_{\displaystyle:=\;\frac{1-\beta}{2\alpha}}
\|s_{k}\|^{2}.
\end{aligned}
\tag{P9}
\]
The equalities in the brackets follow from algebraic cancellation using the admissible stepsize bound $0<\alpha<\frac{2(1-\beta)}{L}$.

\item\textbf{Introduce the Lyapunov function.}  
Recall $\Phi_{k}=f(x_{k})-f_{\inf}+\frac{1-\beta}{2\alpha}\|s_{k}\|^{2}$.  
Subtract $\Phi_{k+1}$ from $\Phi_{k}$ and employ (P9):
\[
\begin{aligned}
\Phi_{k}-\Phi_{k+1}
&= \bigl[f(x_{k})-f(x_{k+1})\bigr]
+\frac{1-\beta}{2\alpha}\bigl(\|s_{k}\|^{2}-\|s_{k+1}\|^{2}\bigr)\\
&\ge \frac{\alpha(1-\beta)}{2}\|\nabla f(x_{k})\|^{2}
\;-\;\frac{1-\beta}{2\alpha}\bigl(\|s_{k+1}\|^{2}-\|s_{k}\|^{2}\bigr).
\end{aligned}
\tag{P10}
\]
But by definition $s_{k+1}=x_{k+1}-x_{k}= \beta s_{k}-\alpha \nabla f(x_{k})$, hence
\[
\|s_{k+1}\|^{2}= \beta^{2}\|s_{k}\|^{2}+\alpha^{2}\|\nabla f(x_{k})\|^{2}
-2\alpha\beta\langle \nabla f(x_{k}),s_{k}\rangle.
\tag{P11}
\]
Using Lemma~\ref{lem:quad} again to bound the inner product term and the admissible range of $\alpha$, one obtains
\[
\|s_{k+1}\|^{2}\le \beta^{2}\|s_{k}\|^{2}+\alpha^{2}\|\nabla f(x_{k})\|^{2}
+\alpha\beta\Bigl(\frac{1-\beta}{\alpha L}\|\nabla f(x_{k})\|^{2}
+\frac{\alpha L}{1-\beta}\|s_{k}\|^{2}\Bigr).
\]
Rearranging gives
\[
\|s_{k+1}\|^{2}-\|s_{k}\|^{2}\le
-\frac{(1-\beta)}{\alpha}\|s_{k}\|^{2}
+\alpha(1-\beta)\|\nabla f(x_{k})\|^{2}.
\tag{P12}
\]
Plugging (P12) into (P10) we finally obtain
\[
\Phi_{k+1}\le \Phi_{k}
-\frac{\alpha(1-\beta)}{2}\|\nabla f(x_{k})\|^{2},
\qquad\text{[Step 1]} \tag{P13}
\]
which is exactly inequality \eqref{1}. \hfill\qedsymbol

\item\textbf{Deriving the rate.}  
Summing (P13) for $k=0,\dots,K-1$ gives
\[
\Phi_{K}\le \Phi_{0}
-\frac{\alpha(1-\beta)}{2}\sum_{k=0}^{K-1}\|\nabla f(x_{k})\|^{2}.
\]
Since $\Phi_{K}\ge 0$, we have
\[
\sum_{k=0}^{K-1}\|\nabla f(x_{k})\|^{2}
\le \frac{2\Phi_{0}}{\alpha(1-\beta)}.
\]
Dividing by $K$ and using $\Phi_{0}\le f(x_{0})-f_{\inf}$ yields the bound (2). \hfill\qedsymbol
\end{enumerate}

\section*{Rate Derivation (With Constants)}
From Theorem~\ref{thm:hb} we have
\[
\min_{0\le t\le K-1}\|\nabla f(x_{t})\|^{2}
\le \frac{2\bigl(f(x_{0})-f_{\inf}\bigr)}{\alpha(1-\beta)K}.
\]
If we set $\alpha = \frac{1-\beta}{L}$ (the largest admissible constant stepsize), the constant becomes
\[
C=\frac{2L\bigl(f(x_{0})-f_{\inf}\bigr)}{(1-\beta)^{2}},
\]
so the rate reads
\[
\min_{0\le t\le K-1}\|\nabla f(x_{t})\|^{2}
\le \frac{C}{K}= \mathcal{O}\!\left(\frac{1}{K}\right).
\]

\section*{Special Cases}
\begin{enumerate}[leftmargin=*]
\item \textbf{Pure Gradient Descent ($\beta=0$).}  
The Lyapunov reduces to $\Phi_{k}=f(x_{k})-f_{\inf}$ and (1) becomes the classic smoothness descent lemma:
\[
f(x_{k+1})\le f(x_{k})-\frac{\alpha}{2}\|\nabla f(x_{k})\|^{2},
\]
recovering the $\mathcal{O}(1/K)$ stationary‑point rate.

\item \textbf{Quadratic Objective.}  
For $f(x)=\frac{1}{2}x^{\top}Qx$ with $Q\succeq \mu I$ and $L$ the largest eigenvalue, the iteration is linear:
\[
\begin{bmatrix}s_{k+1}\\ s_{k}\end{bmatrix}
=
\underbrace{\begin{bmatrix}
(1+\beta)I-\alpha Q & -\beta I\\
I & 0
\end{bmatrix}}_{=:A}
\begin{bmatrix}s_{k}\\ s_{k-1}\end{bmatrix}.
\]
Spectral radius $\rho(A)<1$ is precisely guaranteed by condition \eqref{ass:param}, which explains why the same $\mathcal{O}(1/K)$ bound holds even though the dynamics are oscillatory.

\item \textbf{High Momentum Limit.}  
When $\beta\to 1^{-}$, the admissible stepsize shrinks like $\alpha= \mathcal{O}(1-\beta)$. The product $\alpha(1-\beta)$ stays bounded away from zero, so the rate constant in (2) does not blow up; however, the practical speed‑up diminishes because each step becomes very small.
\end{enumerate}

\end{document}